\section{Conclusion}

The design phase detailed in this chapter marks the critical transition from theoretical requirements to a functional technical solution. By developing a comprehensive design class diagram and mapping it to a robust relational database schema, we have ensured that the \textbf{AgriGov Market} platform is built on a scalable and data-consistent foundation. These architectural choices facilitate the complex logic required for direct trade, allowing for seamless interactions between farmers, buyers, and transporters while providing the Ministry of Agriculture with a reliable interface for market regulation and data monitoring \cite{Kaur2023}.

Furthermore, the presentation of the system's physical architecture and the specific technology stack demonstrates the project's readiness for real-world deployment. The combination of modern development frameworks and a structured deployment strategy ensures that the platform remains performant and accessible across various environments. With the design blueprint finalized and the core interfaces implemented, the technical objectives of this project have been fulfilled. This architectural groundwork provides a stable environment for future enhancements, such as the integration of real-time geolocation or automated analytics, ensuring that the platform can evolve alongside the needs of the Algerian agricultural sector \cite{Belhadi2024}.